
\medskip
\phantomsection\label{lem:sga3-V-7.3}%
\noindent\textbf{Lemma 7.3.}---%
\emph{Retain the notation and hypotheses of
Theorem~\hyperref[thm:sga3-V-7.1]{7.1}. Every point \(z\in X_0\)
that is closed relative to \(S\) has a saturated open neighborhood
\(U_z\) such that the groupoid induced by \(X_\ast\) on \(U_z\) has
a quasi-section.}

The assertion is local on \(S\), so we may suppose that \(S\) is
affine and Noetherian and apply \hyperref[lem:sga3-V-7.2]{the preceding lemma} to the diagram
\SGAThreeDiagram[0.30\textwidth]
  {figures/exp5/p276_lemma73_base_square.png}
  {The diagram to which \hyperref[lem:sga3-V-7.2]{Lemma~7.2} is applied.}
  {fig:sga3-V-7.3-base-square}
  {Polo--Gille re-edition, Expos\'e V}
in \((\mathbf{Sch}/S)\). Thus let \(F\) be a closed subscheme of
\(X_0\) such that
\[
  d_0\bigl(d_0^{-1}(F)\cap d_1^{-1}(z)\bigr)
\]
is finite and nonempty, and such that the restriction of \(d_1\) to
\(d_0^{-1}(F)\) is flat at the points of \(d_1^{-1}(z)\).

Let \(F_1\) and \(F_2\) be the inverse images of \(F\) under \(d_0\)
and \(d_0d'_0=d_0d'_1\), respectively, and let
\(\widetilde d_0,\widetilde d_1,\ldots\) denote the morphisms induced
by \(d_0,d_1,\ldots\). We then have a commutative diagram
\SGAThreeDiagram[0.67\textwidth]
  {figures/exp5/p276_lemma73_F2_diagram.png}
  {The induced diagram on \(F_2,F_1,F\).}
  {fig:sga3-V-7.3-F-diagram}
  {Polo--Gille re-edition, Expos\'e V}
in which the two left-hand squares are Cartesian, the first row is
exact (cf. the diagram
\hyperref[sga3:v:lemma:1-2:equation:0-1-2]{\((0,1,2)\)} in \hyperref[sec:sga3-V-1]{\S1}), and
\(q,\widetilde q\) are the structure morphisms.

We first show that only finitely many points of \(F_1\) lie over
\(z\).%
\footnote{\label{note:sga3-V-39}Editor's note: Details have been added,
and the role of the properness hypothesis on \(d_0\) and \(d_1\) in
Theorem~\hyperref[thm:sga3-V-7.1]{7.1} has been made explicit.
Compare the statement and proof of \hyperref[thm:sga3-V-8.1]{Theorem~8.1}, where this properness
hypothesis is omitted.}
Let \(s\) be the image of \(z\) in \(S\). Since \(F\) is of finite
type over \(S\), the fiber \(\widetilde q^{-1}(s)\) is a Noetherian
scheme. Since \(\widetilde d_0\) is proper,
\[
  \widetilde d_0\bigl(\widetilde d_1^{-1}(z)\bigr)
\]
is a closed subscheme of \(\widetilde q^{-1}(s)\) consisting of
finitely many points. Consequently (cf. EGA~I, 6.2.2), those points
are closed in \(\widetilde q^{-1}(s)\), and also in the fiber
\(q^{-1}(s)\) of \(X_0\), because \(F\) is closed in \(X_0\).

Let \(y\) be one of these points. Since \(q^{-1}(s)\) is of finite
type over \(\kappa(s)\), it has affine open neighborhoods
\[
  \operatorname{Spec}B\quad\text{and}\quad\operatorname{Spec}C
\]
of \(y\) and \(z\), respectively, with \(B\) and \(C\) finitely
generated \(\kappa(s)\)-algebras. The points \(y\) and \(z\)
correspond to maximal ideals
\[
  \mathfrak p\subset B,\qquad \mathfrak q\subset C.
\]
The fields \(B/\mathfrak p\) and \(C/\mathfrak q\) are finite over
\(\kappa(s)\), so
\[
  (B/\mathfrak p)\otimes_{\kappa(s)}(C/\mathfrak q)
\]
is a finite-dimensional \(\kappa(s)\)-algebra. Its maximal ideals
correspond precisely to the points of \(X_0\times_S X_0\) whose
second projection is \(z\) and whose first projection is \(y\).
Thus there are only finitely many points \(u\in X_0\times_S X_0\)
whose second projection is \(z\) and whose first projection belongs
to
\[
  \widetilde d_0\bigl(\widetilde d_1^{-1}(z)\bigr).
\]
Finally, since \(X_1\to X_0\times_S X_0\) has finite fibers, each
such point \(u\) is the image of only finitely many points of \(X_1\).
This proves the assertion.
